

%% ===================================================================
\section{\vrev{방법론적 시사점}}
\label{sec:methodological-implications}
%% ===================================================================

본 연구는 $\DRTO$ 프레임워크의 검증 과정에서 \jrev{네 가지} 방법론적 혁신을
달성하였다. \chg{이를 차례로 서술한다.}


%% -------------------------------------------------------------------
\subsection{\vrev{몬테카를로 시뮬레이션 방법론}}
\label{subsec:mc-methodology}
%% -------------------------------------------------------------------

\chg{첫째, 몬테카를로 시뮬레이션 방법론이다.} $n = 10{,}000$ 표본 $\times$ 4개 조건(C0--C3) $\times$ $10{,}000$개 마스터 시드의
총 $4 \times 10^8$회 $\DRTO$ 계산을 통해, 시뮬레이션 기반 BCM 모델 검증의
새로운 표준을 제시하였다.

\textbf{비중첩 대역 시드(non-overlapping band seeding)} 체계는 마스터 시드
$s = 42$로부터 변수별 시드를 $10{,}000$ 간격의 비중첩 대역\chg{(}$O_{\mathrm{raw}}$:
$[0\text{--}9999]$, $k_O$: $[20000\text{--}29999]$, $U^{(G)}$:
$[30000\text{--}39999]$, $U^{(P)}$: $[40000\text{--}49999]$\chg{)}으로
파생한다. \jrev{H6} 검증에서 6개 변수쌍 모두 $|r| < 0.02$(최대 $|r| = 0.016$)를
충족하여, 시뮬레이션의 완전한 재현성과 변수 간 통계적 독립성이 동시에 보장되었다.

이 체계의 방법론적 의의는 세 가지이다.
\chg{첫째,} 단일 시드로부터 다수의 독립적 난수열을 파생하는 표준화된 규약을 제공한다.
\chg{둘째,} $10{,}000$개 시드에 걸친 $10^8$회 다중 시드 검증(\jrev{H7})에서 Bonferroni 보정
$z$-검정 결과 45개 계수 중 기각이 $0$개로, 특정 시드에 대한 의존성을 완전히
배제한다. \chg{셋째,} 재현성 프로토콜이 명시되어 있어 후속 연구자의 독립적 검증이 가능하다.


%% -------------------------------------------------------------------
\subsection{\jrev{\texorpdfstring{\citet{lee2026drto}}{Lee and Cheung (2026)} $\kappa$ 매개변수의 본 연구 환경 적용 및 \frev{분포 환경 적합성} 시사 방법론}}
\label{subsec:kappa-optimization}
%% -------------------------------------------------------------------

\chg{둘째, $\kappa$ 매개변수의 분포 환경 적합성 시사 방법론이다.} \jrev{본 연구는 곡률 매개변수 $\kappa^{*} = 1.2$를 \chg{\citet{lee2026drto}}에서 다기준 종합점수($f_1$ 실무 유의성,
$f_2$ 비포화율, $f_3$ 효과 크기, $f_4$ 설명력)로 \chg{도출하여 출판된} 매개변수로
받아들이고, 본 연구의 새로운 분포 환경(\chg{C2 가우시안과 C3 파레토})에 그대로 적용하였다. 따라서 $\kappa^{*}$ 도출 자체는 본 연구의
독자적 기여가 아니며, \chg{\citet{lee2026drto}의} 격자 탐색 결과(C1 조건 기준
$F(\kappa^{*}) = 0.9997$)를 \textbf{전제 조건}으로 받아들인다.}

\jrev{본 연구의 방법론적 보완은 다음과 같다. \chg{\citet{lee2026drto}의}
4기준 종합점수 $F(\kappa) = f_1 \cdot f_2 \cdot f_3 \cdot f_4$는
C1 환경 기반으로 정의되어, 새로운 분포 환경(\chg{C2 가우시안과 C3 파레토})에서는
직접 격자 탐색에 부적합하다는 학술적 한계를 지닌다. 본 연구는 격자
탐색의 직접적 대안으로서, 9개 가설(\jrev{H2--H10}) PASS의 \textbf{통합 해석을
통한 \jrev{\frev{분포 환경 적합성}} 종합 명제 도출 방법론}을 제시하였다
(\secref{sec:ch4_hypothesis_summary} \jrev{\frev{분포 환경 적합성}} 종합 명제 참조).
구체적으로, \chg{i) 물리적 안전성(H2), ii) 이론적 일치(H3), iii) 통계적 설명력(H4·H5),
iv) 결과 안정성(H6·H7), v) 메커니즘 발현(H8), vi) 외부 검증(H9·H10)}의 6가지 차원에서 $\kappa = 1.2$의 \jrev{\frev{분포 환경 적합성}}이 집합적으로
\frev{시사된다}.}

\jrev{이러한 \textbf{다중 가설 PASS의 통합 해석} 방법론은 \chg{\citet{lee2026drto}}
매개변수의 새로운 환경 적용성을 \frev{시사}하는 보완적
방법론적 틀로서, 본 연구가 독자적으로 제시하는 기여이다. 본 접근법은
선행연구를 기반으로 분포 환경을 확장하는 후속 연구 일반에 적용 가능하며,
직접적 격자 탐색이 부적합한 환경에서의 매개변수 적합성 검증에 활용될 수 있다.}


%% -------------------------------------------------------------------
\subsection{\vrev{분할 회귀(Split Regression) 방법론}}
\label{subsec:split-regression}
%% -------------------------------------------------------------------

\chg{셋째,} P85.8 기준 분할 회귀는 파레토 분포의 체제 전환(regime change)을
포착하기 위한 방법론적 혁신이다. C3 전체 표본에서 \pdferr{2차 회귀}의
$R^2 = 0.858$에 그치는 잔여 비설명 분산($14.2\%$)을,
P85.8 기준으로 Core($\leq$ P85.8, $n = 8{,}580$)와 Tail($>$ P85.8, $n = 1{,}420$)로
분할함으로써 Core $R^2 = 0.999$, Tail $R^2 = 0.710$으로 개선하였다.

분할 회귀의 방법론적 기여는 다음과 같다.
\chg{i)} 단일 회귀 모형으로 포착할 수 없는 \textbf{구간별 이질적 구조}를
데이터 기반으로 식별하는 체계적 절차를 제시한다.
\chg{ii)} 분할 기준점(P85.8)의 결정이 사전적 이론(CDF 교차)에 근거하므로,
임의적 구간 분할의 자의성 문제를 해소한다.
\chg{iii)} 이 방법론은 BCM 분야에 국한되지 않고, 파레토 분포가 관여하는
금융 리스크 관리, 보험 계리, 자연 재해 모형 등에
범용적으로 확장 가능하다.
